% Image credits for the photographs and AI illustrations of Book 2 (High School).
% Machine-readable license records live in images/book2/CREDITS.md; keep the
% two files in sync. The Book 2 writer owns this file: add one \item per
% photograph (inside the itemize, by chapter) as photographs land.
% Portuguese (Brazilian) edition: translated from frontmatter/image-credits-book2.tex;
% every credit, name, dataset and licence kept exactly.
\chapter*{Créditos das imagens}

Os desenhos, os mapas e os gráficos deste livro são originais; os mapas e
os gráficos desenhados a partir de dados publicados citam a sua fonte na
legenda ou nesta lista. As fotografias são usadas sob as suas respectivas
licenças livres, com agradecimentos aos seus autores; os links completos
das fontes e os endereços das licenças de todas as fotografias estão
listados em \texttt{images/book2/CREDITS.md}, no repositório do livro.

\bigskip
Ao lado das fotografias e dos desenhos originais, algumas figuras são
ilustrações geradas por IA, produzidas com a geração de imagens da OpenAI
a partir de instruções escritas pelos editores do livro, e revisadas
quanto à exatidão científica antes de serem incluídas. Elas mostram
apenas cenas genéricas, nunca um lugar real, uma pessoa real ou uma
amostra real. A instrução completa de cada imagem gerada está registrada
em \texttt{images/book2/ai/PROMPTS.md}, no repositório.

\bigskip
\noindent Fotografias e dados, por capítulo:
\begin{itemize}\small
  \item Cap.~1 (medindo a Terra): a Terra vista da Apollo~17; fotografia da NASA, domínio público (Wikimedia Commons).
  \item Cap.~1: um círculo repetidor, prancha de J.-B. Delambre, \emph{Base du système métrique décimal}, vol.~2 (1807); domínio público (Wikimedia Commons, a partir da BnF Gallica).
  \item Cap.~1: comprimentos dos graus calculados sobre o elipsoide WGS~84 (conjunto de dados EPSG).
  \item Cap.~2 (um planeta habitável): Vênus vista pela Mariner~10 (PIA23791), NASA/JPL-Caltech, domínio público; o monte Sharp visto pelo rover Curiosity (PIA16769), NASA/JPL-Caltech/MSSS, domínio público (Wikimedia Commons).
  \item Cap.~2: dados planetários das fichas técnicas da NASA; curvas de mudança de fase da água das publicações IAPWS SR1-86 e R14-08; limites da zona habitável de Kopparapu e colaboradores (2013); Proxima~b do NASA Exoplanet Archive.
  \item Cap.~3 (o balanço de energia da Terra): poças de degelo sobre o gelo marinho do Ártico (missão ICESCAPE, 2011); fotografia do NASA Goddard Space Flight Center, domínio público (Wikimedia Commons).
  \item Cap.~3: valores do balanço de energia do Sexto Relatório de Avaliação do IPCC, Grupo de Trabalho~I, Figura~7.2; albedos da University of Colorado Boulder (MOSAiC) e do National Snow and Ice Data Center; transporte de calor de Trenberth e Caron (2001); produção primária líquida de Field e colaboradores (1998); dados termoquímicos do NIST Chemistry WebBook.
  \item Cap.~4 (por que as nuvens se formam): cúmulos de bom tempo, fotografia de Ralph F. Kresge, NOAA; a bigorna de um cúmulo-nimbo, NOAA National Severe Storms Laboratory; ambas em domínio público (Wikimedia Commons).
  \item Cap.~4: saturação calculada a partir das publicações IAPWS; tamanhos das gotas do NASA Jet Propulsion Laboratory.
  \item Cap.~5 (trovoadas e ciclones tropicais): o furacão Isabel visto da Estação Espacial Internacional, fotografia de Mike Trenchard, NASA Johnson Space Center; um tornado no Kansas, fotografia de Sean Waugh, NOAA National Severe Storms Laboratory; um raio, fotografia de C.~Clark, NOAA National Severe Storms Laboratory; todas em domínio público (Wikimedia Commons).
  \item Cap.~5: trajetórias de ciclones do IBTrACS versão~4r01 (Gahtan e colaboradores, 2024; Knapp e colaboradores, 2010; NOAA National Centers for Environmental Information); categorias de furacões do NOAA National Hurricane Center; energia dos furacões e condições de formação do NOAA Atlantic Oceanographic and Meteorological Laboratory; frequência dos raios de Christian e colaboradores (2003); dados sobre tornados do NOAA Storm Prediction Center.
  \item Cap.~6 (o oceano em camadas): uma roseta de amostragem, fotografia de Pablo Clemente-Colon, NOAA Photo Library; uma boia perfiladora lançada de um navio de pesquisa, NOAA Global Ocean Monitoring and Observing Program; clorofila oceânica 2002--2010, Jesse Allen e Robert Simmon, NASA Earth Observatory; todas em domínio público (Wikimedia Commons).
  \item Cap.~6: perfis oceânicos do World Ocean Atlas 2023 (Reagan e colaboradores, 2024; NOAA National Centers for Environmental Information); densidade da água do mar da equação de estado internacional (Fofonoff e Millard, 1983); volume do oceano de Charette e Smith (2010); dados sobre as boias do Argo Program Office; ressurgência do NOAA National Ocean Service e de Kessler (2006).
  \item Cap.~7 (o solo, um recurso de formação lenta): um perfil de laterita sobre gnaisse em Kerala, fotografia de Werner Schellmann, CC~BY-SA~2.5; saprólito no centro de Madagascar, fotografia de Diorit, CC~BY-SA~3.0; um perfil de solo, USDA Natural Resources Conservation Service, domínio público (todas Wikimedia Commons).
  \item Cap.~7: taxas de produção de solo e de erosão de Montgomery (2007); fórmulas dos minerais do Mineralogy Database (webmineral.com); espessura do saprólito do USGS Ground Water Atlas of the United States.
  \item Cap.~8 (bacias hidrográficas, rios e cheias): um hidrólogo medindo a vazão de um rio e uma estação fluviométrica, U.S. Geological Survey, domínio público (Wikimedia Commons).
  \item Cap.~8: picos anuais e vazões diárias do rio Potomac perto de Washington, U.S. Geological Survey National Water Information System, estação 01646500; vazão e bacia do Mississippi de Kammerer (1990, USGS); vazão do Amazonas de Callède e colaboradores (2010); segurança diante das águas de cheia do NOAA National Weather Service.
  \item Cap.~9 (os recursos de água doce): o mar de Aral em 2000 e em 2018, NASA Earth Observatory (MODIS/Terra); campos irrigados perto de Sublette, Kansas, NASA (ASTER/Terra), domínio público (Wikimedia Commons).
  \item Cap.~9: estoques mundiais de água segundo Shiklomanov (1993), U.S. Geological Survey Water Science School; o ciclo global da água segundo Trenberth e colaboradores (2007), tal como publicado em Trenberth e Asrar (2014); idades das águas subterrâneas da Núbia de Sturchio e colaboradores (2004); variações do aquífero High Plains de McGuire (2017, USGS); balanço do mar de Aral de Gaybullaev e colaboradores (2012) e Micklin (2007).
  \item Cap.~10 (da deriva continental à expansão do assoalho oceânico): o navio-sonda \emph{JOIDES Resolution}, Ocean Drilling Program, Texas A\&M University, domínio público (Wikimedia Commons).
  \item Cap.~10: a escala de tempo das polaridades de Cande e Kent (1995); os resultados de perfuração da Leg~3 do Deep Sea Drilling Project (1970); a idade do fundo oceânico de Seton e colaboradores (2020), EarthByte, licença CC~BY~4.0.
  \item Cap.~11 (as ondas sísmicas revelam o interior): instalação de um sismômetro temporário em Porto Rico; fotografia de Thomas Pratt, U.S. Geological Survey, domínio público (Wikimedia Commons).
  \item Cap.~11: o Preliminary Reference Earth Model de Dziewonski e Anderson (1981); tempos de percurso empíricos das ondas P e S do Australian Passive Seismic Server (AuSPASS).
  \item Cap.~12 (litosfera e astenosfera): um iceberg, NOAA, domínio público (USGS Water Science School).
  \item Cap.~12: fluxo de calor de Davies e Davies (2010), licença CC~BY~3.0; províncias de fluxo de calor de Roy, Blackwell e Decker (1972); o solidus do manto de Hirschmann (2000); espessura e densidade da crosta de Christensen e Mooney (1995); altitudes do ETOPO1 (NOAA).
  \item Cap.~13 (as placas e os seus movimentos): uma estação de posicionamento por satélite na Cordilheira das Cascatas; fotografia de Dave Sherrod, U.S. Geological Survey, domínio público (Wikimedia Commons). Relevo sombreado do fundo do Pacífico Norte, National Geophysical Data Center, NOAA, domínio público (Wikimedia Commons).
  \item Cap.~13: polos de rotação das placas NNR-MORVEL56 de Argus, Gordon e DeMets (2011); posições das estações do Nevada Geodetic Laboratory (Blewitt, Hammond e Kreemer, 2018); idades dos vulcões de Clague e Dalrymple (1987), U.S. Geological Survey; limites de placas de Bird (2003); continentes do Natural Earth.
  \item Cap.~14 (as dorsais oceânicas): lavas em almofada no Rifte de Galápagos, NOAA Okeanos Explorer Program, domínio público; uma fumarola negra na Dorsal de Juan de Fuca, University of Washington e NOAA, domínio público; gabro, fotografia de Darla Sondrol (GeoDIL), CC0; peridotito em basalto, fotografia de James St.~John, CC~BY~2.0 (todas Wikimedia Commons).
  \item Cap.~14: profundidades do fundo oceânico do ETOPO1 (NOAA) e idades de Seton e colaboradores (2020), EarthByte, licença CC~BY~4.0; leis de profundidade de Parsons e Sclater (1977); o solidus do manto de Hirschmann (2000); polos de rotação NNR-MORVEL56 de Argus, Gordon e DeMets (2011).
  \item Cap.~15 (as zonas de subducção): xisto azul e eclogito, fotografias de James St.~John, CC~BY~2.0; danos do tsunami em Kesennuma, Japão, 2011, fotografia de Dylan McCord, U.S. Navy, domínio público (todas Wikimedia Commons).
  \item Cap.~15: terremotos do catálogo de terremotos do USGS (ComCat); limites de placas e taxas de convergência de Bird (2003); idades do fundo oceânico de Seton e colaboradores (2020), EarthByte, licença CC~BY~4.0; velocidades das placas NNR-MORVEL56 de Argus, Gordon e DeMets (2011); continentes do Natural Earth.
  \item Cap.~16 (os minerais): quartzo, ortoclásio, biotita e olivina, fotografias de Shannon Heinle, Darla Sondrol, Dexter Perkins e Nessa Eull (Geology Digital Image Library), CC0; basalto em lâmina delgada, fotografias de Chmee2, CC~BY-SA~3.0; quartzo de um granito em lâmina delgada, fotografia de Michael C.~Rygel, CC~BY-SA~3.0 (todas Wikimedia Commons).
  \item Cap.~16: a composição da crosta continental de Rudnick e Gao (2014); pesos atômicos padrão da IUPAC/CIAAW; fórmulas e densidades dos minerais de webmineral.com; densidades das rochas graníticas da Sierra Nevada da U.S. Geological Survey (Open-File Report 91-570-A).
  \item Cap.~17 (as rochas ígneas): colunas de basalto da Calçada dos Gigantes, fotografia de Adam Bishop, CC~BY-SA~3.0; El Capitan, Yosemite, U.S. Geological Survey, CC0; basalto, andesito, riolito, diorito, granito, obsidiana, uma rocha porfirítica e duas lâminas delgadas, fotografias de Nessa Eull, Twyla Baker, Joseph H.~Hartman, Mike Davis, Dexter Perkins e Darla Sondrol (Geology Digital Image Library), CC0 (todas Wikimedia Commons).
  \item Cap.~17: resfriamento de diques e soleiras segundo Delaney (1986), U.S. Geological Survey Open-File Report 86-444; composições dos magmas e ordem de cristalização das notas de geologia física de S.~A.~Nelson e de S.~Earle, \emph{Physical Geology} (CC~BY~4.0); a biografia de Bowen da história do Geophysical Laboratory publicada pela Carnegie Science.
  \item Cap.~18 (a geologia do carvão, do petróleo e do gás): plantas fósseis em folhelho rico em carbono, linhito, carvão betuminoso e antracito, fotografias de James St.~John, CC~BY~2.0; uma camada de linhito na Romênia, U.S. Geological Survey, domínio público (todas Wikimedia Commons).
  \item Cap.~18: janela do petróleo, formação do carvão e armadilhas de S.~Earle, \emph{Physical Geology} (CC~BY~4.0); teor de carbono do carvão da U.S. Energy Information Administration; o sistema petrolífero segundo Magoon e Schmoker (U.S. Geological Survey, 2000); as bacias de rifte e de sal do oeste da África Central do U.S. Geological Survey Fact Sheet 2023-3028.
  \item Cap.~19 (a crosta continental): o gnaisse de Acasta, fotografia de Pedroalexandrade, CC~BY-SA~3.0; zircões de Jack Hills, imagens de A.~Cavosie e J.~Valley, domínio público; gnaisses e granitos do Escudo Canadense, fotografia de Chris Light, CC~BY-SA~4.0 (todas Wikimedia Commons).
  \item Cap.~19: composição da crosta de Rudnick e Gao (2014); espessura da crosta de Christensen e Mooney (1995), White, McKenzie e O'Nions (1992) e Karplus e colaboradores (2019); extensão dos continentes de Cogley (1984); idades do fundo oceânico de Seton e colaboradores (2020), EarthByte, licença CC~BY~4.0.
  \item Cap.~20 (a formação das montanhas): o monte Everest visto do ônibus espacial Columbia, NASA, domínio público; camadas de rocha dobradas no desfiladeiro do Kali Gandaki, fotografia de Jean-Marie Hullot, CC~BY~2.0; o cavalgamento de Glarus, fotografia de Christian Heine, CC~BY-SA~3.0; lavas em almofada do ofiolito de Semail, fotografia de Lionel Duchoiselle, CC~BY-SA~3.0 (todas Wikimedia Commons).
  \item Cap.~20: posições passadas da Índia e de Lhasa a partir dos modelos de placas de Müller e colaboradores (2019, 2022), calculadas com o GPlates Web Service (EarthByte); relevo elevado do ETOPO1; terremotos do catálogo do USGS; limites de placas de Bird (2003); continentes do Natural Earth.
  \item Cap.~21 (as rochas metamórficas): clivagem ardosiana ao microscópio, fotografia de Geoloshe, CC0; xisto micáceo com granada de Siros, fotografia de Graeme Churchard, CC~BY~2.0; xisto com cianita e gnaisse de quartzo e biotita, fotografias de James St.~John, CC~BY~2.0 (todas Wikimedia Commons).
  \item Cap.~21: os campos de estabilidade da cianita, da andaluzita e da silimanita de Holdaway e Mukhopadhyay (1993).
  \item Cap.~22 (o magmatismo das zonas de subducção): a coluna eruptiva do monte Santa Helena, 18 de maio de 1980, fotografia de Donald A. Swanson, USGS, domínio público; granodiorito, andesito e a Sierra Nevada, fotografias de James St.~John, CC~BY~2.0 (todas Wikimedia Commons).
  \item Cap.~22: terremotos do catálogo do USGS; vulcões do Smithsonian Global Volcanism Program; limites de placas de Bird (2003); continentes do Natural Earth; o solidus do peridotito seco de Hirschmann (2000), com as densidades do PREM.
  \item Cap.~23 (a destruição das montanhas): os montes Apalaches vistos da Estação Espacial Internacional, fotografia de astronauta ISS033-E-22378, NASA, domínio público; conglomerado de seixos de quartzo, fotografia de James St.~John, CC~BY~2.0 (ambas Wikimedia Commons).
  \item Cap.~23: taxas de denudação de Portenga e Bierman (2011); as cargas sedimentares do Mississippi de Meade e Moody (2010); a constante de decaimento de Ahnert (1970); as Colinas de Varden, a região inventada da série.
  \item Cap.~24 (o calor interno da Terra): a usina geotérmica de Nesjavellir, Islândia, fotografia de Gretar Ívarsson, editada por Fir0002, domínio público (Wikimedia Commons).
  \item Cap.~24: fluxo de calor de Davies e Davies (2010); calor radiogênico da medição de geoneutrinos do KamLAND (2011) e da liberação atual de calor dos quatro isótopos; produção de calor das rochas de Vilà e Fernández (2010).
  \item Cap.~25 (a datação radiométrica e a idade da Terra): William Thomson, lorde Kelvin, 1906, fotografia de Russell and Sons, Londres, Smithsonian Institution Libraries, sem restrições de direitos autorais conhecidas; o meteorito metálico de Canyon Diablo, fotografia de Geoffrey Notkin, CC~BY-SA~2.5 (ambas Wikimedia Commons).
  \item Cap.~25: isótopos do chumbo dos meteoritos de Patterson (1956); a constante de decaimento do rubídio-87 recomendada pela IUGS e pela IUPAC (2022); a história das estimativas a partir de Holmes (1913, 1927), Patterson (1956) e do U.S. Geological Survey; as Colinas de Varden, a região inventada da série.
  \item Cap.~26 (a Terra primitiva e a sua atmosfera): formação ferrífera bandada, Dales Gorge, fotografia de Graeme Churchard, CC~BY~2.0; estromatólitos em Shark Bay, fotografia de Paul Harrison, CC~BY-SA~3.0; camadas vermelhas perto de Sedona, fotografia de Wendy from Pennsylvania, CC~BY-SA~2.0 (todas Wikimedia Commons).
  \item Cap.~26: a luminosidade do Sol jovem de Gough (1981); limites e cores dos éons da Tabela Cronoestratigráfica Internacional.
  \item Cap.~27 (ler os climas do passado): bolhas de ar no gelo antártico, fotografia CSIRO Atmospheric Research, CC~BY~3.0; foraminíferos planctônicos, fotografia de Marci Robinson, USGS, domínio público; grãos de pólen, Dartmouth Electron Microscope Facility, domínio público (todas Wikimedia Commons).
  \item Cap.~27: temperatura de EPICA Dome~C (Jouzel e colaboradores, 2007) e gás carbônico dos testemunhos de gelo antárticos (Bereiter e colaboradores, 2015) via NOAA Paleoclimatology; o empilhamento LR04 (Lisiecki e Raymo, 2005); relevo ETOPO1 (NOAA); NASA GISS (GISTEMP v4); períodos de referência do Sexto Relatório de Avaliação do IPCC, Grupo de Trabalho~I, capítulo~2.
  \item Cap.~28 (a mudança do clima): o Observatório de Mauna Loa; fotografia de Brian Vasel, NOAA, domínio público (Wikimedia Commons).
  \item Cap.~28: Charles David Keeling e John Chin em Mauna Loa, 1997; fotografia de John Miller, NOAA, domínio público (Wikimedia Commons).
  \item Cap.~28: dados do NOAA Global Monitoring Laboratory e da Scripps Institution of Oceanography (gás carbônico), NASA GISS (GISTEMP v4), NOAA Laboratory for Satellite Altimetry (nível do mar), Bereiter e colaboradores (2015) via NOAA Paleoclimatology (testemunhos de gelo), e do Sexto Relatório de Avaliação do IPCC, Grupo de Trabalho~I (forçamentos, atribuição, projeções).
\end{itemize}
